\hypertarget{agent-architectures-must-preserve-scientific-uncertainty}{%
\subsection{Agent architectures must preserve scientific
uncertainty}\label{agent-architectures-must-preserve-scientific-uncertainty}}

\label{exec13:census-agents}

Several systems now make causal analysis a sequence of tool-supported
decisions. Their strongest contribution is often a concrete
representation of the intermediate work: a dataset, a graph, an
estimand, a proof statement or a failed tool call. Comparing these
systems requires examining what each intermediate object means and which
transformations are allowed. A graph stored in agent memory can prevent
a variable-name error while still carrying an unsupported orientation
into every later calculation.

\hypertarget{causal-agent-useful-tool-boundaries-and-unequal-baselines}{%
\subsubsection{Causal Agent: useful tool boundaries and unequal
baselines}\label{causal-agent-useful-tool-boundaries-and-unequal-baselines}}

Han and colleagues store causal graph objects in a dictionary that
agents address by name. This connects independence tests, graph
questions and effect calculations without requiring each operation to
reconstruct the graph from prose. Their CausalTQA examples combine small
generated nonlinear additive Gaussian mechanisms with medical and
marketing vocabulary. The vocabulary tests a language interface; it does
not turn the generated mechanism into a medical or marketing
observation. \citep{exec13_census_causal_agent}

The paper documents informative failures: requesting conditional rather
than marginal independence, passing the wrong variable name, using the
wrong graph operation, and building unnecessarily long tool chains.
Typed variable references and explicit tool signatures are plausible
remedies. The main code-generation baseline, however, executes once,
whereas the agent uses specialized tools and iterative reasoning. Its
measured advantage therefore combines tool design, supplied knowledge
and opportunities for repair. An evaluation intended to isolate
orchestration should give a deterministic workflow and a simpler agent
the same tool implementations and permitted repair budget.

The reported average cost of 3.37 seconds, 2.51 iterations and 5,308
tokens comes from 100 sampled variable, edge and graph questions with a
maximum iteration count of 15. It does not establish the cost of the
full effect workflow, rare failures or deployment under a different
model and platform. An average well below a limit is also insufficient
evidence to remove that limit. The prototype should retain explicit
budgets and report tails, timeouts, failed requests and review effort.
The source's effect-answer tolerance and benchmark oracle require code
inspection before its reported 157 correct answers out of 160 become a
reproduction target. \citep{exec13_census_causal_agent}

\hypertarget{orca-shared-state-without-silently-choosing-a-causal-world}{%
\subsubsection{ORCA: shared state without silently choosing a causal
world}\label{orca-shared-state-without-silently-choosing-a-causal-world}}

ORCA organizes analysis of relational databases around dataset
exploration, discovery and inference. It records table choices, joins,
variable roles, candidate graphs and estimation configurations, with
opportunities for human review. Its REEF benchmark uses a specified
structural causal model inside a relational e-commerce database. This is
directly relevant to a banking proposal: a mistaken join or aggregation
can change the unit of analysis before a causal estimator is called.
\citep{exec13_census_orca}

The important limitation lies between graph discovery and inference. The
paper describes a consensus PAG that retains uncertain endpoints,
followed by a separate ``actionable'\,' DAG. Uncertain directions are
resolved using the highest-ranked candidate, cycle-safe orientation,
reversal or edge removal. Acyclicity does not establish that the
selected orientation is true. An effect computed from this one DAG is
conditional on choices added during that conversion, and cannot
automatically represent all models encoded by the uncertainty-preserving
graph.

For a servicing decision, suppose distress and contact intensity remain
ambiguously related. Directing the edge one way can make distress an
adjustment variable; directing it the other can make it a treatment
descendant. The proposed interface should retain the original graph
object and give every added orientation an explicit assumption record.
It should compare identification across compatible models, return an
effect set where appropriate, or state that a scalar target is not
identified. Ranking an orientation for display is not authorization to
treat it as scientific fact. This is a proposed restriction on reuse,
not behavior shown to exist in ORCA's implementation.

The reported user study also requires careful interpretation. Twelve
participants used ORCA and a baseline with tools, step gates and minimal
state in counterbalanced order. Table 1 gives rounded completion rates
of 0.83 and 0.75. The abstract's 42-percentage-point completion gain is
not consistent with that table; it should not be repeated as a verified
result. Appendix E defines completion as producing an ATE regardless of
accuracy, and discovery success as including both treatment and outcome
nodes. Neither definition establishes correct discovery.
\citep{exec13_census_orca}

The paper's ``CI coverage'\,' description counts an estimated ATE
falling within an interval around a ground-truth effect. That wording
differs from the usual repeated-sampling event in which an estimated
interval contains a fixed true effect. Its exact oracle and
implementation require clarification. Likewise, the full-data test for
the ATE-error condition effect has \(p=.321\), whereas the analysis
excluding large configuration-error outliers reports \(p=.020\) under
OLS. Configuration failures are central outcomes for a workflow product.
They should remain in the primary comparison, with conditional numerical
accuracy reported separately.

\hypertarget{causalpulse-operational-success-has-a-different-denominator}{%
\subsubsection{CausalPulse: operational success has a different
denominator}\label{causalpulse-operational-success-has-a-different-denominator}}

CausalPulse supplies another concrete design: agents coordinate semantic
descriptions, domain rules and causal tools through a utility layer that
includes MCP. Such interfaces can make tools accessible and operations
traceable. The inspected manufacturing study nevertheless distinguishes
successful operation from successful root-cause attribution. Its PSE
example reports Hits@1 of .44, precision .46, recall .40 and F1 .39
despite much higher operational success. Tool availability and
completion are necessary engineering outcomes, but substituting them for
causal correctness would misstate the evidence. \citep{exec13_pulse}

This distinction determines the proposed monitoring interface. It should
record whether the request was parsed, whether the analysis was
admissible, whether the tools completed, whether the scientific target
was answered, and whether the decision was useful under the declared
loss. These events need separate denominators. A scientific refusal can
be a correct answer, while a complete report can be wrong. Banking
acceptance would additionally need appropriate target-bank evidence; the
public prototype can already test these distinctions with generated
known-truth cases.

\hypertarget{causalsmith-a-checked-proof-still-needs-a-checked-statement}{%
\subsubsection{CausalSmith: a checked proof still needs a checked
statement}\label{causalsmith-a-checked-proof-still-needs-a-checked-statement}}

CausalSmith combines a Lean library with proposal, formalization, proof
and presentation stages. Its paper reports 8,179 declarations and a
catalog of 144 runs: 14 accepted, 59 downgraded and 71 failed. Those
counts include reruns as separate entries and exclude a legacy track.
They describe one campaign with changing model versions, not independent
draws from a fixed research-task population. The authors explicitly
acknowledge that research quality ratings lack independent human
validation and that the error rate of their LLM statement-match audit is
unknown. \citep{exec13_census_causalsmith}

The source's most useful distinction is between checking a proof against
a formal statement and checking whether that statement expresses the
intended claim. A compiled artifact can assume a hard step as an axiom,
weaken a conclusion, strengthen an assumption or make a predicate
vacuous. The paper describes mechanical scans and dependency checks
alongside language-model statement review. These are complementary
checks with different assurances. The latter does not acquire a formal
guarantee merely by sitting next to a proof checker.

For a proposed bank-effect theorem, the statement record should expose
the unit, intervention, estimand, time ordering, conditioning
information, positivity and dependence assumptions. The proof
environment should bind the toolchain, imported definitions and axioms.
A separate review should compare the English claim with that exact
formal statement and test whether its assumptions are satisfiable in the
intended model. A theorem conditional on observed confounders cannot be
presented as solving unmeasured confounding in operational data. Formal
correctness, semantic fidelity and scientific relevance each need their
own evidence.

The near-term reuse choice is therefore selective. Retain typed graph
memory, explicit intermediate state, preserved failures, and formal
statement/dependency records as design candidates. Defer executable
adoption of these newly inspected systems until exact code,
dependencies, licenses and relevant numerical or proof paths have been
reviewed. This permits a concrete prototype architecture now without
claiming that paper descriptions have already been integrated or
independently reproduced.

\label{exec13:census-agents-end}
